\exercisegroup

\begin{enumerate}
\def\labelenumi{\arabic{enumi}.}
\item
  一个图 \(\mathbf{G}\) 是集合 \(\mathbf{E}\) 上等价关系的图，当且仅当 \(\mathrm{pr}_1\mathbf{G} = \mathbf{E}\) , \(\mathbf{G} \circ \overline{\mathbf{G}} \circ \mathbf{G} = \mathbf{G}\) ，以及 \(\Delta_{\mathbf{E}} \subset \mathbf{G}\) ( \(\Delta_{\mathbf{E}}\) 为 \(\mathbf{E}\) 的对角线).
\item
  如果 \(\mathbf{G}\) 是一个图，使得
\end{enumerate}

\[
\mathbf {G} \circ \overline {{{{\mathbf {G}}}}} ^ {- 1} \circ \mathbf {G} = \mathbf {G},
\]

证明 \(\overline{\mathbf{G}}\circ \mathbf{G}\) 并且 \(\mathbf{G}\circ \overline{\mathbf{G}}\) 分别是 \(\mathfrak{pr}_1\mathbf{G}\) 并且 \(\mathfrak{pr}_2\mathbf{G}\) 上的等价关系的图。

\begin{enumerate}
\def\labelenumi{\arabic{enumi}.}
\setcounter{enumi}{2}
\item
  设E是一个集合，A是E的一个子集，R是与映射 \(X \rightarrow X \cap A\) （从 \(\mathfrak{P}(E)\) 到 \(\mathfrak{P}(E)\) ）相关联的等价关系. 证明存在一个从 \(\mathfrak{P}(A)\) 到商集上 \(\mathfrak{P}(E)/R\) 的双射.
\item
  让 \(G\) 是集合上一个等价关系的图 \(E\) 证明：如果 \(A\) 是一个图，使得 \(A \subset G\) 并且 \(\operatorname{pr}_1 A = E\) （分别） \(\operatorname{pr}_2 A = E\) ），则
\end{enumerate}

\[
\mathbf {G} \circ \mathbf {A} = \mathbf {G} (\text { 分别 } \quad \mathbf {A} \circ \mathbf {G} = \mathbf {G});
\]

此外，如果 B 是任意图，我们有 \((G \cap B) \circ A = G \cap (B \circ A)\) （分别） \(A \circ (G \cap B) = G \cap (A \circ B)\) ).

\begin{enumerate}
\def\labelenumi{\arabic{enumi}.}
\setcounter{enumi}{4}
\item
  证明：集合 E 上等价关系的图的任意交集仍是 E 上某个等价关系的图。给出集合 E 上两个等价关系的例子，使得它们的图的并集不是 E 上某个等价关系的图。
\item
  设 G, H 是 E 上两个等价的图。则 \(G \circ H\) 是 E 上等价的图当且仅当 \(G \circ H = H \circ G\) 。图 \(G \circ H\) 此时是 E 上包含 G 和 H 的所有等价图的交集。
\item
  让 \(G_{0}\) , \(G_{1}\) , \(H_{0}\) , \(H_{1}\) 是集合 E 上四个等价的图，使得 \(G_{1} \cap H_{0} = G_{0} \cap H_{1}\) 并且 \(G_{1} \circ H_{0} = G_{0} \circ H_{1}\) 对于每一个 \(x \in E\) ，让 \(R_{0}\) （分别） \(S_{0}\) ) 是由等价关系 \(\mathbf{G}_{1}(x)\) （分别） \(H_{1}(x)\) ) 在 \((x, y) \in \mathbf{G}_{0}\) （分别） \((x, y) \in \mathbf{H}_{0}\) 上诱导的关系。证明存在商集 \(\mathbf{G}_{1}(x)/\mathbf{R}_{0}\) 到商集上
\end{enumerate}

\[
\mathrm{H} _ {1} (x) / \mathrm{S} _ {0}.
\]

的一个双射（若 \(\mathbf{A} = \mathbf{G}_1(x) \cap \mathbf{H}_1(x)\) ，证明两个商集都与 \(\mathbf{A}\) 关于由 \(\mathbf{R}_0\) 在 \(\mathbf{A}\) 上诱导的等价关系所成的商集一一对应；该关系等价于由 \(\mathbf{S}_0\) 在 \(\mathbf{A}\) 上诱导的关系 .)

\begin{enumerate}
\def\labelenumi{\arabic{enumi}.}
\setcounter{enumi}{7}
\item
  设E、F为两个集合，设R为F上的一个等价关系，并设f为从E到F的一个映射。若S为在f下R的逆像所构成的等价关系，且若 \(A = f \langle E \rangle\) ，定义 E/S 到上 \(A/R_{A}\) 的一个典范双射.
\item
  设 F、G 为两个集合，R 为 F 上的一个等价关系，p 为 F 到 F/R 的典范映射，且 f 为 G 到 F/R 的一个满射。证明存在一个集合 E、一个 E 到 F 的满射 g，以及一个 E 到 G 的满射 h，使得 \(p \circ g = f \circ h\) .
\item
  \begin{enumerate}
  \def\labelenumii{(\alph{enumii})}
  \tightlist
  \item
    如果 \(R\{x, y\}\) 若任意关系，则“ \(R\{x, y\}\) 并且 \(R\{y, x\}\) ”是一个对称关系。在什么条件下它在集合 E 上是自反的？
  \end{enumerate}
\end{enumerate}

*(b) 设 \(\mathbb{R}\{x, y\}\) 设为一个集合上的自反且对称的关系 \(\mathbf{E}\) . 设 \(\mathbb{S}\{x, y\}\) 为关系“存在整数 \(n > 0\) 以及一个序列 \((x_i)_{0 \leqslant i \leqslant n}\) 的元素 \(\mathbf{E}\) 使得 \(x_0 = x\) , \(x_n = y\) , 并且对每个指标 \(i\) 使得 \(0 \leqslant i < n\) , \(\mathbb{R}\{x_i, x_{i+1}\}\) ”。证明 \(\mathbb{S}\{x, y\}\) 是上的一个等价关系 \(\mathbf{E}\) ，并且其图是 \(\mathbf{E}\) 上所有包含 \(\mathbb{R}\) 的图的等价关系的图中最小的。关于 \(\mathbb{S}\) 被称为连通分量 \(\mathbf{E}\) 关于该关系 \(\mathbb{R}\) .

\begin{enumerate}
\def\labelenumi{(\alph{enumi})}
\setcounter{enumi}{2}
\tightlist
\item
  让 \(\mathfrak{F}\) 设 A 为 E 的子集所构成的集合，使得对于每一对元素 \((y,z)\) 使得 \(y\in A\) 并且 \(z\in E - A\) ，我们有“非 R \(\{y,z\} "\) .
\end{enumerate}

对于每一个 \(x \in \mathbf{E}\) ，证明这些集合的交集 \(\mathbf{A} \in \mathfrak{F}\) 使得 \(x \in \mathbf{A}\) 是...的连通分量 \(x\) 关于该关系 \(\mathbf{R}_{*}\) .

\begin{enumerate}
\def\labelenumi{\arabic{enumi}.}
\setcounter{enumi}{10}
\tightlist
\item
  \begin{enumerate}
  \def\labelenumii{(\alph{enumii})}
  \tightlist
  \item
    设 \(\mathbb{R}\{x, y\}\) 是集合 \(\mathbf{E}\) 上的一个自反且对称的关系。 \(\mathbf{R}\) 被称为1阶非传递的，如果对于任意四个不同的元素 \(x, y, z, t\) （都属于 \(\mathbf{E}\) ），这些关系 \(\mathbb{R}\{x, y\}, \mathbb{R}\{x, z\}, \mathbb{R}\{x, t\}, \mathbb{R}\{y, z\}\) 以及 \(\mathbb{R}\{y, t\}\) 蕴含 \(\mathbb{R}\{z, t\}\) 。 \(\mathbf{A}\) （ \(\mathbf{E}\) 的子集）称为关于关系 \(\mathbb{R}\) 稳定的，如果 \(\mathbb{R}\{x, y\}\) 对所有 \(x\) 与 \(y\) 属于 \(\mathbf{A}\) 都成立。如果 \(a\) 与 \(b\) 是 \(\mathbf{E}\) 的两个不同元素，使得 \(\mathbb{R}\{a, b\}\) ，证明集合 \(\mathbf{C}(a, b)\) ，即由元素 \(x \in \mathbf{E}\) 满足 \(\mathbb{R}\{a, x\}\) 与 \(\mathbb{R}\{b, x\}\) 者组成的集合，是稳定的，并且
  \end{enumerate}
\end{enumerate}

\[
\mathbf {C} (x, y) = \mathbf {C} (a, b)
\]

对于每一对不同元素 \(x, y\) 的 \(C(a, b)\) 集合 \(C(a, b)\) （对于每个有序对 \((a, b)\) 使得 \(R\{a, b\}\) ）以及关于 \(R\) 的连通分量（练习10）中由单个元素组成的那些分量称为 \(E\) 关于该关系 \(R\) 的构成部分。证明 \(E\) 的两个不同构成部分的交集至多包含一个元素，并且如果 \(A, B, C\) 是三个两两不同的构成部分，则集合 \(A \cap B, B \cap C, C \cap A\) 中至少有一个是空集。

\begin{enumerate}
\def\labelenumi{(\alph{enumi})}
\setcounter{enumi}{1}
\item
  反之，设 \((\mathbf{X}_{\lambda})_{\lambda \in \mathbf{L}}\) 是集合的一个覆盖， \(\mathbf{E}\) 由非空子集组成， \(\mathbf{E}\) 并且具有以下性质：（1）如果 \(\lambda, \mu\) 是两个不同的指标， \(\mathbf{X}_{\lambda} \cap \mathbf{X}_{\mu}\) 至多包含一个元素；（2）若 \(\lambda, \mu, \nu\) 是三个不同的指标，则这三个集合中至少有一个 \(\mathbf{X}_{\lambda} \cap \mathbf{X}_{\mu}, \mathbf{X}_{\mu} \cap \mathbf{X}_{\nu}, \mathbf{X}_{\nu} \cap \mathbf{X}_{\lambda}\) 为空。设 \(\mathbf{R}\{x, y\}\) 为关系“存在 \(\lambda \in \mathbf{L}\) 使得 \(x \in \mathbf{X}_{\lambda}\) 并且 \(y \in \mathbf{X}_{\lambda}\) ；证明 \(\mathbf{R}\) 在 \(\mathbf{E}\) 上是自反的、对称的且为1阶非传递的，并且 \(\mathbf{X}_{\lambda}\) 是 \(\mathbf{E}\) 关于 \(\mathbf{R}\) .
\item
  的组成部分。* 类似地，一个在E上自反且对称的关系 \(R\{x, y\}\) 被称为n-3阶非传递的，如果对于E中任意一族 \((x_{i})_{1 \leqslant i \leqslant n}\) 不同的元素，关系 \(R\{x_{i}, x_{j}\}\) 对于每一对 \((i, j) \neq (n-1, n)\) 蕴含 \(R\{x_{n-1}, x_{n}\}\) 成立。将(a)和(b)的结果推广到任意阶的非传递关系。证明一个p阶非传递的关系对于所有q也是q阶非传递的。 \textgreater{} 第～*页
\end{enumerate}
